\chapter{Data Analysis as Statistical Learning}
\label{chap:data-analysis-statistical-learning}

\section{Why This Course Begins with Statistical Learning}

Data analysis is not just the act of computing summaries or fitting models.  It is the disciplined process of using data to learn something about a system, make predictions about future observations, and communicate conclusions in a way that can be checked, reproduced, and defended.  In this course, we use the phrase \emph{statistical learning} to emphasize that our models are tools for learning from data rather than formulas that are true simply because we wrote them down.

This course serves as a bridge.  It connects the probability, statistics, and mathematical tools developed in earlier operations analysis courses to the modeling and machine-learning methods that appear later.  The immediate goal is to build the statistical modeling intuition needed for regression, exploratory data analysis, model validation, and diagnostics.  The longer-term goal is to prepare you to use more advanced methods without treating them as black boxes.

A recurring theme is the statement, often summarized as ``all models are wrong.''  The point is not that models are useless.  The point is that a model is a deliberate simplification.  It keeps the structure we think is useful and ignores or absorbs the rest as noise, approximation error, measurement error, or unexplained variation.  A good analyst understands both sides: what the model is trying to capture and what it is failing to capture.

\section{What Is Data?}

A data set is a collection of measurements, records, or observations.  Each observation corresponds to one unit of analysis: a person, ship, aircraft, market, mission, experiment, time step, or other object being studied.  Each variable records one attribute of that observation.

In statistical modeling, we usually distinguish between \emph{input variables} and an \emph{output variable}.  Input variables are the measured quantities we use to explain or predict something.  They are also called \emph{predictors}, \emph{features}, \emph{covariates}, or \emph{independent variables}.  The output variable is the quantity we are trying to explain or predict.  It is also called the \emph{response} or \emph{dependent variable}.

For example, suppose we want to model a service member's physical fitness test score.  Possible inputs might include age, training frequency, prior score, sleep, and injury status.  The output could be the next test score.  In a logistics setting, inputs might be demand history, location, distance, and inventory status, while the output could be delivery time or stockout risk.  In a classification problem, inputs might be sensor readings and the output might be whether a track is classified as friendly, hostile, or unknown.

We collect the input variables for a single observation into a vector.  For an observation with $p$ inputs, write
\[
  \mathbf{x} = (x_1,x_2,\ldots,x_p)^T \in \mathbb{R}^p.
\]
Here $p$ is the number of input variables, and $x_j$ is the value of the $j$th input for this observation.  The corresponding output is denoted by $Y$ or, for a realized observed value, by $y$.

\paragraph{Math Foundations Connection.}
Math Foundations treats a vector as an ordered collection of quantities that belong to the same object, not merely as an arrow in physical space.  That is exactly how we use $\mathbf{x}$ here.  The vector $\mathbf{x}$ is one observation's collection of input measurements.  Its entries may be physical quantities, demographic attributes, encoded categories, time summaries, or engineered features.  The geometric interpretation of vectors will become useful later, but the first interpretation is simpler: a feature vector is a structured record of measured information.  See the Math Foundations textbook, Chapter 2.

\section{A First Statistical Model}

The central modeling idea in this chapter is that an output variable $Y$ is related to input variables $\mathbf{X} = (X_1,X_2,\ldots,X_p)^T$ through an unknown relationship plus random error:
\begin{equation}
  Y = f(\mathbf{X}) + \epsilon.
  \label{eq:general-stat-model}
\end{equation}
The function $f$ represents the systematic part of the relationship between the inputs and the output.  The random variable $\epsilon$ represents the part of $Y$ not explained by $f(\mathbf{X})$.

A simple example with two predictors is
\begin{equation}
  Y = \beta_0 + \beta_1 x_1 + \beta_2 x_2 + \epsilon.
  \label{eq:first-linear-example}
\end{equation}
In this model, $\beta_0$ is an intercept, $\beta_1$ and $\beta_2$ are coefficients, and $\epsilon$ is the unobserved error.  The inputs $x_1$ and $x_2$ are observed.  The output $Y$ is observed after the fact.  The error term $\epsilon$ is not directly observed; it is the difference between the observed output and the systematic part of the model.

Equation~\eqref{eq:general-stat-model} is intentionally general.  The function $f$ might be linear, nonlinear, simple, complicated, interpretable, or opaque.  In this course, we often begin with linear models because they are mathematically transparent and remain important for inference.  But the broader statistical learning framework is not limited to linear models.

We usually impose basic assumptions on the error term.  A common starting point is
\begin{align}
  \mathbb{E}[\epsilon] &= 0, \label{eq:error-mean-zero}\\
  \operatorname{Var}(\epsilon) &> 0, \label{eq:error-positive-variance}\\
  \epsilon &\text{ is independent of } \mathbf{X}. \label{eq:error-independent}
\end{align}
The mean-zero assumption says that the error does not systematically push the output upward or downward once the systematic relationship has been accounted for.  The positive-variance assumption says that there is genuine unexplained variability.  The independence assumption says that the leftover error is not still carrying predictable information from the inputs.  These assumptions are not decorations.  They are part of the model's claim about the world, and later chapters will show how violations of assumptions affect inference, prediction, and diagnostics.

Since $f$ is unknown, we estimate it using data.  We denote the estimated function by $\hat{f}$.  Given a new input vector $\mathbf{x}$, the corresponding predicted output is
\begin{equation}
  \widehat{Y} = \hat{f}(\mathbf{x}).
  \label{eq:prediction-general}
\end{equation}
The hat notation signals that the quantity is estimated or predicted rather than directly observed or known.

\paragraph{Math Foundations Connection.}
When the model is linear in the input variables, the systematic part can be written compactly as
\[
  f(\mathbf{x}) = \beta_0 + \beta_1 x_1 + \cdots + \beta_p x_p
  = \beta_0 + \boldsymbol{\beta}^T \mathbf{x},
\]
where $\boldsymbol{\beta} = (\beta_1,\ldots,\beta_p)^T$.  The term $\boldsymbol{\beta}^T\mathbf{x}$ is an inner product.  Equivalently, it is a linear combination of the entries of $\mathbf{x}$, with coefficients $\beta_1,\ldots,\beta_p$.  This is why linear algebra becomes the natural language of regression: the model combines measured inputs by weighting them and adding them together.  See the Math Foundations textbook, Chapters 2, 3, and 5.

\section{Inference and Prediction}

There are two major reasons to estimate $f$: inference and prediction.  They are related, but they are not the same task.

\subsection{Inference}

\emph{Inference} is about understanding the relationship between the inputs and the output.  In the linear model
\[
  Y = \beta_0 + \beta_1 x_1 + \beta_2 x_2 + \epsilon,
\]
inference focuses on questions such as:
\begin{itemize}
  \item What is the meaning of $\beta_1$ and $\beta_2$?
  \item Are these coefficients positive, negative, or plausibly zero?
  \item Which predictors are most relevant to $Y$?
  \item How much does the output change when one input changes and the others are held fixed?
  \item What assumptions are required before we trust those conclusions?
\end{itemize}

Inference is about learning structure.  If a commander asks why readiness changed, which drivers matter, or whether a particular intervention had an effect, then the analysis is inferential.  Accuracy still matters, but the model must also be interpretable enough to support an explanation.

\subsection{Prediction}

\emph{Prediction} is about estimating the output for new or future inputs.  Given a new vector $\mathbf{x}_{\text{new}}$, we use the fitted model to compute
\[
  \widehat{Y}_{\text{new}} = \hat{f}(\mathbf{x}_{\text{new}}).
\]
Prediction focuses on questions such as:
\begin{itemize}
  \item Given these inputs, what output should we expect?
  \item How accurate is the prediction likely to be?
  \item How much uncertainty remains even after we use the model?
  \item How well will the model perform on cases it has not yet seen?
\end{itemize}

Prediction is about learning what is likely to happen next.  If an analyst is building a model to forecast demand, estimate casualty risk, classify a sensor return, or predict equipment failure, then the model may be judged primarily by its predictive performance.

\subsection{The Same Model Can Serve Both Goals}

The same mathematical model can sometimes support both inference and prediction, but the priorities differ.  A highly interpretable model may be preferred for inference even if a more flexible model predicts slightly better.  Conversely, in a pure prediction task, a less interpretable model may be acceptable if it performs substantially better on new data.  A defensible analysis states which goal is primary and evaluates the model accordingly.

\section{Regression and Classification}

A first way to categorize statistical learning problems is by the type of output variable.

In a \emph{regression} problem, the output is numerical or quantitative.  Examples include predicting fuel consumption, time to complete a task, exam score, cost, weight, or sales.  A regression model returns a number.  A simple linear regression model has the form
\begin{equation}
  \widehat{y} = wx + b,
  \label{eq:linear-regression-first}
\end{equation}
where $w$ is a slope or weight and $b$ is an intercept or bias term.

In a \emph{classification} problem, the output is a category or class label.  Examples include predicting pass/fail, disease/no disease, hostile/friendly/unknown, or low/medium/high risk.  A classifier may return a class directly, but many classifiers first estimate probabilities for each class and then choose the most likely class.

For binary classification, where $Y\in\{0,1\}$, logistic regression models the probability of class 1 as
\begin{equation}
  \Pr(Y=1\mid \mathbf{x}) = \sigma(w_1x_1+w_2x_2+b),
  \label{eq:logistic-first}
\end{equation}
where
\begin{equation}
  \sigma(t) = \frac{1}{1+e^{-t}}
\end{equation}
is the logistic or sigmoid function.  The expression $w_1x_1+w_2x_2+b$ is a score.  The sigmoid maps that score to a probability between 0 and 1.

For multiclass classification, a model may compute several scores, often called \emph{logits}.  For example, an image classifier might produce a cat logit, dog logit, and ship logit.  Those logits can then be transformed into class probabilities.  The key idea is that classification models often separate two steps: first produce class-specific scores, then convert those scores into a decision or probability distribution.

\CoreCompress{
\section{Linear and Nonlinear Models}

The previous examples introduce a second distinction: whether the model class is linear or nonlinear.  This distinction is separate from the regression/classification distinction.  We can have linear regression, nonlinear regression, linear classification, and nonlinear classification.  For now, "linear" means the familiar geometric picture: lines, planes, and hyperplanes in the predictors as written.  Chapter~\ref{chap:linear-models-beyond-straight-lines} later sharpens this by showing that a model can still be linear after transforming predictors, as long as it remains linear in the unknown coefficients.

A \emph{linear regression model} predicts a numerical output using a linear function of the inputs:
\[
  \widehat{y}=\beta_0 + \beta_1x_1 + \cdots + \beta_p x_p.
\]
With one predictor, this is a line.  With two predictors, it is a plane.  With more predictors, it is a hyperplane.

A \emph{linear classification model} uses a linear score to separate classes.  In two dimensions, the decision boundary is a line.  Logistic regression is a standard example: the model computes a linear score and then converts that score into a probability.

But data are not always linear.  Sometimes the pattern in a regression problem is curved rather than straight.  For example, a one-dimensional nonlinear regression model might be written as
\begin{equation}
  \widehat{y} = w(x-a)^2 + b.
  \label{eq:nonlinear-regression-first}
\end{equation}
This model can represent a curved relationship because the input appears through the transformed term $(x-a)^2$.

Similarly, a nonlinear classifier may use a nonlinear score before applying the sigmoid function.  For example,
\begin{equation}
  \Pr(Y=1\mid \mathbf{x}) = \sigma\left(\gamma_0 + \gamma_1(x_1-h)^2 + \gamma_2(x_2-k)^2\right).
  \label{eq:nonlinear-classifier-first}
\end{equation}
This model can create a curved decision boundary in the $(x_1,x_2)$ plane.  The boundary is no longer forced to be a straight line.

The important point is that \emph{data shape} and \emph{model shape} must be considered together.  If the true relationship is nonlinear but the model is linear, the model may miss essential structure.  If the model is extremely flexible, it may capture real structure but also chase random noise.  Later chapters develop this tradeoff carefully.

\subsection{Neural Networks as Repeated Linear and Nonlinear Operations}

A fully connected neural network layer is another example of this linear/nonlinear pattern.  A basic layer takes an input vector $\mathbf{x}$, forms a vector of linear combinations, and then applies a nonlinear activation function:
\begin{equation}
  \mathbf{h} = \sigma(W\mathbf{x}+\mathbf{b}).
  \label{eq:nn-layer-first}
\end{equation}
Here $W$ is a matrix of weights, $\mathbf{b}$ is a vector of intercepts or biases, and $\sigma$ is applied component-by-component.  Without $\sigma$, stacked layers are still only affine transformations (linear maps plus shifts).  The nonlinearity is what lets neural networks represent more complicated functions.

This course does not begin with neural networks because the basic modeling questions are easier to see in linear regression.  However, the same foundational ideas remain: inputs, outputs, parameters, loss, assumptions, prediction error, and the need to justify the modeling choices.
}{
\section{Linear and Nonlinear Models}

Models can be either linear or nonlinear, and that distinction is separate from whether the task is regression or classification.  A linear regression model predicts a numerical output using a linear function of the inputs,
\[
  \widehat{y}=\beta_0+\beta_1x_1+\cdots+\beta_px_p.
\]
With one predictor this is a line; with two predictors it is a plane; with more predictors it is a hyperplane.  A linear classifier uses a similar linear score to separate or score classes, often followed by a probability transformation such as the logistic function.

Nonlinear models allow more flexible shapes.  For example, a regression model might include a term such as $(x-a)^2$, and a classifier might use a curved score before converting it to a probability.  The key modeling question is whether the model class can capture the structure in the data without chasing noise.  Later, Chapter~\ref{chap:linear-models-beyond-straight-lines} sharpens the word ``linear'' by showing that transformed predictors can still lead to a model that is linear in its unknown coefficients.

\subsection{Neural Networks as Repeated Linear and Nonlinear Operations}

A fully connected neural-network layer has the form
\begin{equation}
  \mathbf{h}=\sigma(W\mathbf{x}+\mathbf{b}).
\end{equation}
Here $W$ is a weight matrix, $\mathbf{b}$ is a bias vector, and $\sigma$ is applied componentwise.  Without $\sigma$, stacked layers are still only affine transformations: linear maps plus shifts.  The nonlinearity is what lets neural networks represent more complicated functions.
}
\CoreCompress{
\section{Statistical Concepts Used Throughout the Course}

Before developing prediction error and model diagnostics, we need a small set of statistical concepts.  This section reviews expectation, variance, covariance, correlation, and bias.  These ideas will reappear throughout the course.

\subsection{Expectation}

The expectation of a random variable is its probability-weighted average.  For a discrete random variable $X$ with possible values $x_1,\ldots,x_m$ and probabilities $p_1,\ldots,p_m$, the expectation is
\begin{equation}
  \mathbb{E}[X] = \sum_{i=1}^m p_i x_i.
  \label{eq:expectation-discrete}
\end{equation}
The probabilities weight the possible outcomes.  Outcomes that are more likely contribute more to the average.

\paragraph{Math Foundations Connection.}
If we collect the possible outcomes into a vector $\mathbf{x}=(x_1,\ldots,x_m)^T$ and the probabilities into a vector $\mathbf{p}=(p_1,\ldots,p_m)^T$, then
\[
  \mathbb{E}[X] = \mathbf{p}^T\mathbf{x}.
\]
Thus expectation can be viewed as an inner product between a probability vector and an outcome vector.  The probability vector must have nonnegative entries that sum to one.  This connects the statistical idea of a weighted average to the linear algebra idea of an inner product.  See the Math Foundations textbook, Chapters 2, 3, and 5.

Expectation has two especially important algebraic properties.  First, for constants $a$ and $b$,
\begin{equation}
  \mathbb{E}[aX+b] = a\mathbb{E}[X]+b.
  \label{eq:affine-expectation}
\end{equation}
Second, for random variables $X$ and $Y$,
\begin{equation}
  \mathbb{E}[X+Y] = \mathbb{E}[X] + \mathbb{E}[Y].
  \label{eq:linearity-expectation}
\end{equation}
This second property does \emph{not} require $X$ and $Y$ to be independent.  This is why we call it linearity of expectation.

For the affine rule, the reasoning is straightforward:
\[
  \mathbb{E}[aX+b]
  = \mathbb{E}[aX] + \mathbb{E}[b]
  = a\mathbb{E}[X] + b,
\]
since a constant has expectation equal to itself.

\subsection{Variance}

Variance measures the average squared deviation from the mean.  If $\mu_X = \mathbb{E}[X]$, then
\begin{equation}
  \operatorname{Var}(X)=\mathbb{E}\left[(X-\mu_X)^2\right].
  \label{eq:variance-definition}
\end{equation}
A useful computational identity is
\begin{equation}
  \operatorname{Var}(X)=\mathbb{E}[X^2]-\left(\mathbb{E}[X]\right)^2.
  \label{eq:variance-shortcut}
\end{equation}
To see why, expand the square:
\begin{align*}
  \operatorname{Var}(X)
  &= \mathbb{E}\left[(X-\mu_X)^2\right] \\
  &= \mathbb{E}\left[X^2 - 2\mu_X X + \mu_X^2\right] \\
  &= \mathbb{E}[X^2] - 2\mu_X\mathbb{E}[X] + \mu_X^2 \\
  &= \mathbb{E}[X^2] - 2\mu_X^2 + \mu_X^2 \\
  &= \mathbb{E}[X^2] - \mu_X^2.
\end{align*}
Since $\mu_X=\mathbb{E}[X]$, this gives Equation~\eqref{eq:variance-shortcut}.

Variance changes predictably when we scale and shift a random variable.  For constants $a$ and $b$,
\begin{equation}
  \operatorname{Var}(aX+b)=a^2\operatorname{Var}(X).
  \label{eq:variance-affine}
\end{equation}
The constant $b$ does not affect variance because shifting all outcomes by the same amount does not change their spread.  The scale factor $a$ multiplies deviations by $a$, so squared deviations are multiplied by $a^2$.

\subsection{Covariance and Correlation}

Variance measures how one random variable varies around its mean.  Covariance measures how two random variables vary together.  If $\mu_X=\mathbb{E}[X]$ and $\mu_Y=\mathbb{E}[Y]$, then
\begin{equation}
  \operatorname{Cov}(X,Y)
  = \mathbb{E}\left[(X-\mu_X)(Y-\mu_Y)\right].
  \label{eq:covariance-definition}
\end{equation}
If $X$ tends to be above its mean when $Y$ is above its mean, the covariance tends to be positive.  If $X$ tends to be above its mean when $Y$ is below its mean, the covariance tends to be negative.  If their deviations do not move together in a systematic linear way, the covariance may be near zero.

Covariance appears in the variance of a sum:
\begin{equation}
  \operatorname{Var}(X+Y)
  = \operatorname{Var}(X)+\operatorname{Var}(Y)+2\operatorname{Cov}(X,Y).
  \label{eq:variance-sum}
\end{equation}
A direct derivation comes from writing centered variables.  Let $\mu_X=\mathbb{E}[X]$ and $\mu_Y=\mathbb{E}[Y]$.  Then
\begin{align*}
  \operatorname{Var}(X+Y)
  &= \mathbb{E}\left[\left((X+Y)-(\mu_X+\mu_Y)\right)^2\right] \\
  &= \mathbb{E}\left[\left((X-\mu_X)+(Y-\mu_Y)\right)^2\right] \\
  &= \mathbb{E}\left[(X-\mu_X)^2\right]
   + \mathbb{E}\left[(Y-\mu_Y)^2\right]
   + 2\mathbb{E}\left[(X-\mu_X)(Y-\mu_Y)\right] \\
  &= \operatorname{Var}(X)+\operatorname{Var}(Y)+2\operatorname{Cov}(X,Y).
\end{align*}

If $X$ and $Y$ are independent, then $\operatorname{Cov}(X,Y)=0$.  The converse is not generally true: zero covariance does not necessarily imply independence.  Covariance detects linear co-movement, while dependence can be nonlinear.

Correlation is a scaled version of covariance:
\begin{equation}
  \rho_{XY}
  = \operatorname{Corr}(X,Y)
  = \frac{\operatorname{Cov}(X,Y)}{\sqrt{\operatorname{Var}(X)}\sqrt{\operatorname{Var}(Y)}}.
  \label{eq:correlation-definition}
\end{equation}
Correlation is unitless and lies between $-1$ and $1$ when both variances are positive.  The scaling by the standard deviations makes correlation easier to compare across contexts than covariance.

\paragraph{Math Foundations Connection.}
For observed data, sample correlation can be interpreted geometrically.  If we center two data vectors by subtracting their sample means, then their sample correlation is the cosine of the angle between the centered vectors.  This is why correlation measures linear alignment: two centered data vectors pointing in nearly the same direction have correlation near $1$; two pointing in opposite directions have correlation near $-1$; and two nearly orthogonal centered vectors have correlation near $0$.  See the Math Foundations textbook, Chapters 2, 3, and 5.

\subsection{Bias}

An estimator is a rule for using data to estimate an unknown quantity.  If the unknown parameter is $\theta$ and the estimator is $\hat{\theta}$, then the bias of the estimator is
\begin{equation}
  \operatorname{Bias}(\hat{\theta})
  = \mathbb{E}[\hat{\theta}] - \theta.
  \label{eq:bias-definition}
\end{equation}
Bias is systematic error.  It asks whether the estimator tends to miss in a particular direction over repeated samples.  Bias is not the same thing as random error.  An estimator can be unbiased and still vary substantially from sample to sample.  Conversely, a biased estimator may have low variability but systematically target the wrong value.

Bias can be induced by measurement choices, data collection procedures, model assumptions, or the deliberate use of a simplified model.  This is one reason the phrase ``all models are wrong'' matters.  A model can be useful while still introducing bias.  The analyst's job is not to pretend the model is perfect; it is to understand what kind of error the model introduces and whether that error is acceptable for the task.
}{
\section{Statistical Concepts Used Throughout the Course}

This course repeatedly uses a small set of statistical ideas: expectation, variance, covariance, correlation, and bias.  The full reference text gives more derivation detail; the Core Reader keeps the definitions and interpretations needed for the chapters that follow.

\subsection{Expectation}

For a discrete random variable $X$ with possible values $x_1,\ldots,x_m$ and probabilities $p_1,\ldots,p_m$,
\begin{equation}
  \mathbb{E}[X]=\sum_{i=1}^m p_i x_i.
\end{equation}
Expectation is a probability-weighted average.  In vector form, if $\mathbf{x}$ stores the outcomes and $\mathbf{p}$ stores the probabilities, then $\mathbb{E}[X]=\mathbf{p}^T\mathbf{x}$.  This is a weighted dot product.  See the Math Foundations textbook, Chapters 2, 3, and 5.

Two properties matter throughout the course:
\begin{equation}
  \mathbb{E}[aX+b]=a\mathbb{E}[X]+b,
  \qquad
  \mathbb{E}[X+Y]=\mathbb{E}[X]+\mathbb{E}[Y].
\end{equation}
Linearity of expectation does not require independence.

\subsection{Variance}

Variance measures average squared deviation from the mean:
\begin{equation}
  \operatorname{Var}(X)=\mathbb{E}\left[(X-\mu_X)^2\right],
  \qquad \mu_X=\mathbb{E}[X].
\end{equation}
A useful equivalent form is
\begin{equation}
  \operatorname{Var}(X)=\mathbb{E}[X^2]-\left(\mathbb{E}[X]\right)^2.
\end{equation}
Shifting a random variable does not change its variance, while scaling by $a$ multiplies variance by $a^2$.

\subsection{Covariance and Correlation}

Covariance measures how two random variables move together:
\begin{equation}
  \operatorname{Cov}(X,Y)=\mathbb{E}\left[(X-\mu_X)(Y-\mu_Y)\right].
\end{equation}
Correlation rescales covariance to a unitless number:
\begin{equation}
  \rho_{XY}=\frac{\operatorname{Cov}(X,Y)}{\sqrt{\operatorname{Var}(X)}\sqrt{\operatorname{Var}(Y)}}.
\end{equation}
For observed data, correlation can be interpreted as linear alignment between centered data vectors.  This connects directly to dot products, norms, and angles; see the Math Foundations textbook, Chapters 2, 3, and 5.

\subsection{Bias}

If $\theta$ is an unknown parameter and $\widehat{\theta}$ is an estimator, then
\begin{equation}
  \operatorname{Bias}(\widehat{\theta})=\mathbb{E}[\widehat{\theta}]-\theta.
\end{equation}
Bias is systematic error.  An estimator can be unbiased but variable, or biased but stable.  Later chapters use this distinction to discuss prediction error and model flexibility.
}
\section{Reproducibility and Defensibility}

A data analysis must be more than numerically correct at one moment on one computer.  It should also be reproducible and defensible.

\subsection{Reproducibility}

Reproducibility means that the same code applied to the same data produces the same results.  A reproducible workflow has several features:
\begin{itemize}
  \item the data source is clear;
  \item preprocessing steps are documented;
  \item transformations are recorded;
  \item modeling choices are traceable;
  \item assumptions and exclusions are stated;
  \item results can be rerun without guessing what happened.
\end{itemize}

A useful standard is this: your future self, or another analyst, should be able to rerun the analysis and understand the major decisions without needing a private explanation from you.

\subsection{Defensibility}

Defensibility concerns what we can explain to a customer, commander, decision maker, or technical reviewer.  A defensible analysis answers questions such as:
\begin{itemize}
  \item Why did we choose this model?
  \item What assumptions does the method rely on?
  \item What alternatives did we consider?
  \item What are the limitations of the data?
  \item How much uncertainty remains?
  \item What conclusion is supported, and what conclusion would overstate the evidence?
\end{itemize}

Reproducibility is about what we did.  Defensibility is about what we can justify.  Both are essential.  A result can be reproducible but not defensible if it repeatedly applies the wrong method.  A result can sound defensible but not be reproducible if the workflow cannot be checked.  Good data analysis requires both.

\CoreCompress{
\section{Math Foundations Source Map}

This source map records the Math Foundations support used in this chapter and where those ideas appear in the completed Math Foundations textbook.

\begin{center}
\begin{tabular}{|p{0.24\textwidth}|p{0.36\textwidth}|p{0.28\textwidth}|}
\hline
\textbf{Chapter 1 use} & \textbf{Math Foundations connection} & \textbf{MF textbook location} \\
\hline
Feature vectors for observations & $\mathbb{R}^n$ vector notation, vector entries, subvectors, and vectors as collections of quantities belonging to the same object & Ch02 Secs. 2.1, 2.2, 2.6 \\
\hline
Linear models as weighted inputs & Scalar multiplication and linear combinations of vectors & Ch03 Secs. 3.4, 3.5 \\
\hline
Inner-product form $\boldsymbol{\beta}^T\mathbf{x}$ & Dot product definition, dot product properties, and examples of weighted sums & Ch03 Sec. 3.5; Ch05 Secs. 5.1, 5.3 \\
\hline
Neural-network layers as repeated linear operations plus nonlinearities & Matrix-vector multiplication as row inner products and as linear combinations of columns & Ch03 Sec. 3.5; Ch05 Secs. 5.1, 5.3; Ch07 Secs. 7.1, 7.6.1, 7.6.2 \\
\hline
Correlation as linear alignment & Dot products, Euclidean norm, distance, angle, and orthogonality & Ch05 Secs. 5.1, 5.3, 5.5, 5.7, 5.8 \\
\hline
Reasoning from assumptions to conclusions & Logic, truth tables, implications, biconditionals, and De Morgan's laws & Ch01 Secs. 1.1, 1.7 \\
\hline
\end{tabular}
\end{center}
}{
\section{Math Foundations Source Map}

The Core Reader keeps the short Math Foundations connections where they clarify the chapter.  The full Math Foundations textbook-location map for this chapter appears in the full reference textbook.
}
\section{Chapter Summary}

This chapter introduced data analysis as statistical learning.  A data set consists of observations and variables.  In supervised learning, we distinguish input variables from an output or response variable.  We model their relationship using
\[
  Y = f(\mathbf{X}) + \epsilon,
\]
where $f$ is the systematic relationship and $\epsilon$ is the unexplained error.

The two major modeling goals are inference and prediction.  Inference asks what the model says about the relationship between inputs and output.  Prediction asks how accurately the model can estimate outputs for new inputs.  Regression problems involve quantitative outputs, while classification problems involve categorical outputs.  Models may be linear or nonlinear, and modern methods such as neural networks build complex functions by combining linear operations with nonlinear activations.

Finally, the chapter reviewed expectation, variance, covariance, correlation, and bias.  These concepts provide the mathematical language needed to discuss model error, uncertainty, assumptions, and the limits of prediction.  The next chapter uses this foundation to study prediction error, model flexibility, overfitting, and generalization.
